\documentclass{beamer}
\usetheme{Madrid}
\usepackage{amsmath}
\usepackage{graphicx}
\title{Variational Quantum Simulation of 2D Spin Models in Qiskit}
\date{}
\begin{document}

\begin{frame}{Introduction to Quantum Spin Models}
\begin{itemize}
\item Spin lattice models such as the Ising and Heisenberg models are fundamental in condensed matter physics, describing magnetism and correlated quantum many-body systems [oai_citation:0‡en.wikipedia.org](https://en.wikipedia.org/wiki/Transverse-field_Ising_model#:~:text=Image%3A%20%7B%5Cdisplaystyle%20H%3D,j%7D%5Cright) [oai_citation:1‡nature.com](https://www.nature.com/articles/s41467-024-46402-9?error=cookies_not_supported&code=e0d86c6c-d400-4a7e-9dda-2e6669633ed5#:~:text=intrinsic%20gates%20to%20design%20the,in%20a%20magnetic%20field%20B).
\item These models have non-commuting terms (e.g. $Z_iZ_j$ vs $X_i$) that make classical simulation intractable for large lattices [oai_citation:2‡en.wikipedia.org](https://en.wikipedia.org/wiki/Transverse-field_Ising_model#:~:text=Image%3A%20%7B%5Cdisplaystyle%20H%3D,j%7D%5Cright).
\item Quantum computers can efficiently encode the $2^N$-dimensional Hilbert space of $N$ spins; variational algorithms like VQE are promising for finding ground states on near-term (NISQ) devices [oai_citation:3‡learning.quantum.ibm.com](https://learning.quantum.ibm.com/tutorial/variational-quantum-eigensolver#:~:text=Variational%20quantum%20algorithms%20are%20promising,Given%20that).
\end{itemize}
\end{frame}

\begin{frame}{2D Transverse-Field Ising Model}
\begin{itemize}
\item The 2D transverse-field Ising Hamiltonian on a square lattice is given by
\[
H_{\text{Ising}} = -J \sum_{\langle i,j\rangle} Z_i Z_j \;-\; h \sum_{i} X_i,
\]
where $\langle i,j\rangle$ denotes nearest neighbors, $J$ is the coupling, and $h$ the transverse field strength [oai_citation:4‡en.wikipedia.org](https://en.wikipedia.org/wiki/Transverse-field_Ising_model#:~:text=Image%3A%20%7B%5Cdisplaystyle%20H%3D,j%7D%5Cright).
\item The first term ($Z_iZ_j$) favors aligned or anti-aligned spins on adjacent sites; the second term ($X_i$) applies a perpendicular magnetic field.
\item Here $X_i,Z_i$ are Pauli matrices on qubit (spin) $i$. The competition of these non-commuting terms leads to a quantum phase transition in large lattices.
\end{itemize}
\end{frame}

\begin{frame}{2D Isotropic Heisenberg Model}
\begin{itemize}
\item The 2D isotropic Heisenberg model on a square lattice is
\[
H_{\text{Heis}} = J \sum_{\langle i,j\rangle} \bigl(X_i X_j + Y_i Y_j + Z_i Z_j\bigr),
\]
with the same coupling $J$ for $XX$, $YY$, and $ZZ$ interactions.
\item This SU(2)-symmetric model describes exchange interactions between neighboring spins. It is a prototypical model for quantum magnetism.
\item On a 2×2 lattice (4 spins), each spin has 2 neighbors (in an open boundary scenario, there are 4 edges in total).
\end{itemize}
\end{frame}

\begin{frame}{Variational Quantum Eigensolver (VQE)}
\begin{itemize}
\item VQE is a hybrid quantum-classical algorithm for finding the ground state energy of a Hamiltonian [oai_citation:5‡learning.quantum.ibm.com](https://learning.quantum.ibm.com/tutorial/variational-quantum-eigensolver#:~:text=Variational%20quantum%20algorithms%20are%20promising,Given%20that).
\item A parameterized quantum circuit $U(\boldsymbol{\theta})$ (the ansatz) prepares a trial state $\ket{\psi(\boldsymbol{\theta})} = U(\boldsymbol{\theta})\ket{0}$; the cost function is the energy expectation
\[
E(\boldsymbol{\theta}) = \bra{0}\,U^\dagger(\boldsymbol{\theta})\,H\,U(\boldsymbol{\theta})\ket{0},
\]
which is iteratively minimized [oai_citation:6‡nature.com](https://www.nature.com/articles/s41467-024-46402-9?error=cookies_not_supported&code=e0d86c6c-d400-4a7e-9dda-2e6669633ed5#:~:text=The%20first%20variational%20algorithm%20that,qubit%20photonic%20processor).
\item Classical optimizers update the parameters $\boldsymbol{\theta}$ to drive $E(\boldsymbol{\theta})$ toward the ground-state energy. VQE requires fewer coherent gates than quantum phase estimation, making it suitable for NISQ devices [oai_citation:7‡learning.quantum.ibm.com](https://learning.quantum.ibm.com/tutorial/variational-quantum-eigensolver#:~:text=Variational%20quantum%20algorithms%20are%20promising,Given%20that) [oai_citation:8‡nature.com](https://www.nature.com/articles/s41467-024-46402-9?error=cookies_not_supported&code=e0d86c6c-d400-4a7e-9dda-2e6669633ed5#:~:text=In%20comparison%20to%20QPE%2C%20which,optimization%20of%20the%20variational%20parameters).
\end{itemize}
\end{frame}

\begin{frame}{Hamiltonian Construction in Qiskit}
\begin{itemize}
\item In Qiskit, Hamiltonians are expressed as sums of Pauli operators using tensor products [oai_citation:9‡medium.com](https://medium.com/qiskit/introducing-qiskit-algorithms-with-qiskit-runtime-primitives-d89703ecfca3#:~:text=hamiltonian%20%3D%20SparsePauliOp.from_list%28%5B%28%E2%80%9CII%E2%80%9D%2C%20,0.01%29%2C%20%28%E2%80%9CYX%E2%80%9D%2C%200.1).
\item One can construct a \texttt{SparsePauliOp} (or \texttt{PauliSumOp}) from a list of Pauli strings and coefficients. For example, a 2×2 Ising Hamiltonian uses terms like \verb|"ZZII"|-J, \verb|"IZZI"|-J, \verb|"XIII"|-h, etc.
\item Pseudo-code example:
\begin{verbatim}
from qiskit.quantum_info import SparsePauliOp
# Define Pauli terms for H_Ising on 4 qubits
paulis = [("ZZII", -J), ("IZZI", -J), ("ZIZI", -J), ("IIZZ", -J),
         ("XIII", -h), ("IXII", -h), ("IIXI", -h), ("IIIX", -h)]
H = SparsePauliOp.from_list(paulis)
\end{verbatim}
\item This leverages Qiskit's operator algebra to build $H$ directly from Pauli products [oai_citation:10‡medium.com](https://medium.com/qiskit/introducing-qiskit-algorithms-with-qiskit-runtime-primitives-d89703ecfca3#:~:text=hamiltonian%20%3D%20SparsePauliOp.from_list%28%5B%28%E2%80%9CII%E2%80%9D%2C%20,0.01%29%2C%20%28%E2%80%9CYX%E2%80%9D%2C%200.1).
\end{itemize}
\end{frame}

\begin{frame}{Ansatz Design and Optimization}
\begin{itemize}
\item We use the \texttt{TwoLocal} ansatz from Qiskit: alternating single-qubit rotation layers and entangling gates [oai_citation:11‡medium.com](https://medium.com/qiskit/introducing-qiskit-algorithms-with-qiskit-runtime-primitives-d89703ecfca3#:~:text=estimator%20%3D%20Estimator,%E2%80%9Cry%E2%80%9D%2C%20%E2%80%9Crz%E2%80%9D%5D%2C%20entanglement_blocks%3D%E2%80%9Dcz%E2%80%9D).
\item For instance, use rotations around Y and Z axes (\texttt{ry}, \texttt{rz}) on each qubit, and controlled-Z (\texttt{cz}) entangling gates. Layers (reps) can be increased for expressiveness.
\item In a 2×2 lattice, entangling gates are applied along the edges (e.g. CZ between qubits 0–1, 0–2, 1–3, 2–3) to reflect the lattice geometry.
\item The VQE is executed on a statevector simulator (e.g.\ \texttt{BasicAer.get_backend('statevector_simulator')}) with a classical optimizer (COBYLA, SLSQP, etc.) to minimize the energy.
\end{itemize}
\end{frame}

\begin{frame}[fragile]{Implementation: 2×2 Transverse Ising}
\begin{itemize}
\item Construct the 4-qubit Hamiltonian and run VQE as follows:
\begin{verbatim}
from qiskit.quantum_info import SparsePauliOp
from qiskit.circuit.library import TwoLocal
from qiskit.algorithms import VQE
from qiskit.algorithms.optimizers import COBYLA
from qiskit import BasicAer

# Define Pauli terms for H = -J sum Z_i Z_j - h sum X_i
paulis = [
("ZZII",-J), ("IZZI",-J), ("ZIZI",-J), ("IIZZ",-J),
("XIII",-h), ("IXII",-h), ("IIXI",-h), ("IIIX",-h)
]
H = SparsePauliOp.from_list(paulis)

ansatz = TwoLocal(num_qubits=4, rotation_blocks=['ry','rz'],
                 entanglement_blocks='cz', reps=2)

result = VQE(ansatz, COBYLA(),
            quantum_instance=BasicAer.get_backend('statevector_simulator')
           ).compute_minimum_eigenvalue(operator=H)
E0 = result.eigenvalue
\end{verbatim}
\item This yields the estimated ground-state energy $E_0$ of the 2×2 Ising model. The exact ground energy (from diagonalizing $H$) can be used to verify accuracy.
\end{itemize}
\end{frame}

\begin{frame}[fragile]{Implementation: 2×2 Heisenberg}
\begin{itemize}
\item Construct the Hamiltonian with $XX$, $YY$, and $ZZ$ terms for each neighbor:
\begin{verbatim}
edges = [(0,1),(0,2),(1,3),(2,3)]
paulis = []
for (i,j) in edges:
   pX = ["I"]*4; pX[i] = "X"; pX[j] = "X"
   paulis.append(("".join(pX), J))
   pY = ["I"]*4; pY[i] = "Y"; pY[j] = "Y"
   paulis.append(("".join(pY), J))
   pZ = ["I"]*4; pZ[i] = "Z"; pZ[j] = "Z"
   paulis.append(("".join(pZ), J))
H = SparsePauliOp.from_list(paulis)

ansatz = TwoLocal(num_qubits=4, rotation_blocks=['ry','rz'],
                 entanglement_blocks='cz', reps=2)

result = VQE(ansatz, COBYLA(),
            quantum_instance=BasicAer.get_backend('statevector_simulator')
           ).compute_minimum_eigenvalue(operator=H)
E0 = result.eigenvalue
\end{verbatim}
\item This provides the ground-state energy estimate for the 2×2 Heisenberg model. The true ground energy (via exact diagonalization) can similarly validate VQE performance.
\end{itemize}
\end{frame}

\begin{frame}{Ground State Energy vs Coupling}
\begin{itemize}
\item By varying the coupling parameter ($h$ for the Ising model, $J$ for the Heisenberg model) and rerunning VQE (or exact diagonalization), one obtains the ground-state energy $E_0$ as a function of that parameter.
\item For example, for a fixed $J=1$ in the Ising model, compute $E_0(h)$ for $h$ in some range (e.g.\ $0$ to $2$). Similarly, compute $E_0(J)$ for the Heisenberg model.
\item These data points can be plotted to show the variation of $E_0$ with the coupling (see figure caption below). The resulting curves (monotonic in these simple cases) serve as a benchmark of the model’s physics and the VQE’s accuracy.
\item Such plots are useful to observe trends or phase transitions in the ground state as parameters change, and to compare VQE vs exact results.
\end{itemize}
\end{frame}

\begin{frame}{Lattice Connectivity and Circuit Illustration}
\begin{itemize}
\item The 2×2 lattice connectivity forms a square: qubits 0–1–3–2 (with 0–2 and 1–3 as vertical edges, 0–1 and 2–3 as horizontal edges).
\item In the quantum circuit, each qubit corresponds to a spin. Entangling gates (e.g.\ CZ or CNOT) are placed between qubits that are neighbors on the lattice to capture interactions.
\item For instance, in the ansatz we apply CZ gates on pairs (0,1), (0,2), (1,3), (2,3) to reflect the 2D nearest-neighbor edges.
\item (A circuit diagram would show single-qubit rotations on each qubit and CZ gates along the edges of a 2×2 grid, with the pattern of entanglement matching the lattice graph.)
\end{itemize}
\end{frame}

\begin{frame}{Scaling and Limitations}
\begin{itemize}
\item The Hilbert space dimension grows as $2^N$ for $N$ spins, so exact simulation is intractable for large $N$. VQE helps but still faces challenges as system size grows.
\item The number of variational parameters and required measurements increases with lattice size. Large parameter counts can lead to optimization difficulties (such as barren plateaus where gradients vanish) [oai_citation:12‡nature.com](https://www.nature.com/articles/s41467-024-46402-9?error=cookies_not_supported&code=e0d86c6c-d400-4a7e-9dda-2e6669633ed5#:~:text=Another%20field%20of%20research%20concerns,New%20ways%20to).
\item NISQ limitations (finite coherence time, gate errors) restrict circuit depth and qubit connectivity. Shallow ansätze may not capture the true ground state entanglement; deeper circuits are more accurate but more error-prone.
\item To date, most VQE demonstrations on spin models have been limited to small lattices (few qubits) [oai_citation:13‡nature.com](https://www.nature.com/articles/s41467-024-46402-9?error=cookies_not_supported&code=e0d86c6c-d400-4a7e-9dda-2e6669633ed5#:~:text=intrinsic%20gates%20to%20design%20the,in%20a%20magnetic%20field%20B). Scaling to larger 2D systems is an ongoing research challenge.
\end{itemize}
\end{frame}

\begin{frame}[allowframebreaks]{References}
\begin{itemize}
\item IBM Quantum (2021), \emph{Variational Quantum Eigensolver} tutorial [oai_citation:14‡learning.quantum.ibm.com](https://learning.quantum.ibm.com/tutorial/variational-quantum-eigensolver#:~:text=Variational%20quantum%20algorithms%20are%20promising,Given%20that) [oai_citation:15‡learning.quantum.ibm.com](https://learning.quantum.ibm.com/tutorial/variational-quantum-eigensolver#:~:text=Like%20many%20classical%20optimization%20problems%2C,the%20Qiskit%20Runtime%20Estimator%20directly).
\item Wright, K. \emph{et al.}, Qiskit blog “Introducing Qiskit Algorithms with Primitives” (2023) [oai_citation:16‡medium.com](https://medium.com/qiskit/introducing-qiskit-algorithms-with-qiskit-runtime-primitives-d89703ecfca3#:~:text=hamiltonian%20%3D%20SparsePauliOp.from_list%28%5B%28%E2%80%9CII%E2%80%9D%2C%20,0.01%29%2C%20%28%E2%80%9CYX%E2%80%9D%2C%200.1) [oai_citation:17‡medium.com](https://medium.com/qiskit/introducing-qiskit-algorithms-with-qiskit-runtime-primitives-d89703ecfca3#:~:text=estimator%20%3D%20Estimator,%E2%80%9Cry%E2%80%9D%2C%20%E2%80%9Crz%E2%80%9D%5D%2C%20entanglement_blocks%3D%E2%80%9Dcz%E2%80%9D).
\item Gras, A. \emph{et al.}, “Quantum many-body simulations on digital quantum computers”, \emph{Nat. Commun.} \textbf{15}, 104 (2024) [oai_citation:18‡nature.com](https://www.nature.com/articles/s41467-024-46402-9?error=cookies_not_supported&code=e0d86c6c-d400-4a7e-9dda-2e6669633ed5#:~:text=Another%20field%20of%20research%20concerns,New%20ways%20to) [oai_citation:19‡nature.com](https://www.nature.com/articles/s41467-024-46402-9?error=cookies_not_supported&code=e0d86c6c-d400-4a7e-9dda-2e6669633ed5#:~:text=The%20first%20variational%20algorithm%20that,qubit%20photonic%20processor).
\item Wikipedia (2025), “Transverse-field Ising model” (accessed June 2025) [oai_citation:20‡en.wikipedia.org](https://en.wikipedia.org/wiki/Transverse-field_Ising_model#:~:text=Image%3A%20%7B%5Cdisplaystyle%20H%3D,j%7D%5Cright).
\end{itemize}
\end{frame}

\end{document}
